% ---
% title: Rating migration
% tags: finance, multiplication-rule, total-probability, bayes
% ---

\begin{problem}
	\begin{context}
		A \emph{credit rating} (for example from S\&P, Moody's or Fitch) measures
		how likely a bond issuer is to fail to pay. Bonds rated BBB$-$ or better
		are \emph{investment grade} (IG); lower-rated bonds are \emph{high yield}
		(HY), as they must pay higher interest (a higher yield) to attract
		investors. Ratings
		change over time; the yearly probabilities of moving from one rating to
		another are collected in a \emph{rating migration (transition) matrix}.
	\end{context}
	The one-year rating transitions of a bond are the following. An
	investment-grade (IG) bond stays IG with probability $0.93$, becomes high
	yield (HY) with probability $0.06$ and defaults (D) with probability $0.01$.
	An HY bond becomes IG with probability $0.05$, stays HY with probability
	$0.85$ and defaults with probability $0.10$. Default is final: a defaulted
	bond stays in default. A bond is IG today, and in each year the transition
	follows these probabilities, depending only on the rating at the start of
	that year.
	\begin{enumerate}
		\item[(a)] What is the probability that the bond defaults within 2 years?
		\item[(b)] Given that it defaulted within 2 years, what is the probability
		      that it was first downgraded to HY?
	\end{enumerate}
\end{problem}

\begin{solution}
	\idea Draw the two-year tree. The probability of a path is the product of
	the probabilities along it (multiplication rule). In (a) we add the paths
	that end in D (total probability). In (b) we divide the ``via HY'' path by
	the total of (a) (Bayes' formula).

	\step{The tree.} The numbers given are conditional probabilities: for
	example, $0.10$ is the probability of default in a year \emph{given} that
	the bond is HY at the start of that year. Write $X_1$ for the rating after
	one year and $X_2$ for the rating after two years. The tree starts from
	IG today. Each path gets the product of its branch probabilities:
	\begin{center}
		\begin{tikzpicture}[
				grow=right, level distance=30mm,
				level 1/.style={sibling distance=30mm},
				level 2/.style={sibling distance=11mm},
				every node/.style={font=\small},
				edge from parent/.style={draw,gray,->},
				lab/.style={font=\scriptsize,fill=white,inner sep=1pt,text=black,midway},
				st/.style={circle,draw=navy,fill=navy!8,inner sep=1pt,minimum size=17pt},
				dd/.style={circle,draw=gold!80!black,fill=gold!25,inner sep=1pt,minimum size=17pt}]
			\node[st] {IG}
			child {node[dd] {D}
					child {node[dd] (l7) {D}
							edge from parent node[lab] {$1$}}
					edge from parent node[lab] {$0.01$}}
			child {node[st] {HY}
					child {node[dd] (l6) {D}
							edge from parent node[lab] {$0.10$}}
					child {node[st] (l5) {HY}
							edge from parent node[lab] {$0.85$}}
					child {node[st] (l4) {IG}
							edge from parent node[lab] {$0.05$}}
					edge from parent node[lab] {$0.06$}}
			child {node[st] {IG}
					child {node[dd] (l3) {D}
							edge from parent node[lab] {$0.01$}}
					child {node[st] (l2) {HY}
							edge from parent node[lab] {$0.06$}}
					child {node[st] (l1) {IG}
							edge from parent node[lab] {$0.93$}}
					edge from parent node[lab] {$0.93$}};
			\foreach \nn/\pp in {l1/{0.93\cdot0.93=0.8649},
					l2/{0.93\cdot0.06=0.0558},
					l3/{0.93\cdot0.01=0.0093},
					l4/{0.06\cdot0.05=0.0030},
					l5/{0.06\cdot0.85=0.0510},
					l6/{0.06\cdot0.10=0.0060},
					l7/{0.01\cdot1=0.0100}}
			\node[anchor=west,xshift=2mm,font=\scriptsize] at (\nn.east) {$\pp$};
			\node[font=\scriptsize,gray] at (0,4.8) {today};
			\node[font=\scriptsize,gray] at (3,4.8) {after 1 year};
			\node[font=\scriptsize,gray] at (6,4.8) {after 2 years};
		\end{tikzpicture}
	\end{center}
	As a check, the seven path probabilities add up to $1$.

	\step{(a)} The bond defaults within 2 years exactly when its path ends in
	D. There are three such paths (gold nodes), and they are mutually
	exclusive (a bond follows only one path), so we add their probabilities:
	\begin{itemize}
		\item default in the first year: $0.01$ (after that it stays in D with
		      probability $1$);
		\item IG after one year, then default:
		      $P(X_1=\text{IG})\,P(X_2=\text{D}\mid X_1=\text{IG})=0.93\cdot0.01=0.0093$;
		\item HY after one year, then default:
		      $P(X_1=\text{HY})\,P(X_2=\text{D}\mid X_1=\text{HY})=0.06\cdot0.10=0.006$.
	\end{itemize}
	Each product is the multiplication rule $P(EF)=P(E)\,P(F\mid E)$. So
	\[ P(X_2=\text{D}) = 0.01+0.0093+0.006 = 0.0253. \]
	This is the law of total probability, with the rating after one year
	(IG, HY or D) as the partition.

	\step{(b)} Let $E=\{X_2=\text{D}\}$ (default within 2 years) and
	$H=\{X_1=\text{HY}\}$ (first downgraded to HY). We need $P(H\mid E)$. The
	event $HE$ is the single path IG$\to$HY$\to$D, with probability $0.006$.
	By the definition of conditional probability (this is Bayes' formula):
	\[ P(H\mid E)=\frac{P(HE)}{P(E)}=\frac{0.06\cdot0.10}{0.0253}
		=\frac{0.006}{0.0253}\approx 0.237. \]
	In words: among the bonds that default within two years, about $23.7\%$
	went through HY first.

	\begin{answer}
		(a) $0.0253$ \qquad (b) $\dfrac{0.006}{0.0253}\approx 0.237$
	\end{answer}
	\begin{remark}
		The other two default paths split the rest:
		$0.01/0.0253\approx0.395$ of the defaults happen in the first year, and
		$0.0093/0.0253\approx0.368$ happen in the second year after the bond
		stayed IG. The three shares add up to $1$.
		Notice that the HY path is the least likely of the three, although HY
		bonds default ten times as often as IG bonds: only $6\%$ of the bonds
		reach HY at all. \textbf{Common mistake:} answering $0.10$ in (b), which
		is $P(E\mid H)$, not $P(H\mid E)$.
	\end{remark}
\end{solution}
